\chapter{Implementation}
\label{c:implementation}

\noindent
This chapter documents the entire analytical pipeline, from prompt construction to the statistical
procedures that underpin each hypothesis test. The aim is reproducibility, so every
transformation applied to the data is explained, along with the decision made wherever more than
one treatment was defensible.

The pipeline has five stages. Prompts are constructed and models are queried
(Section~\ref{impl:prompts}); human annotations are gathered under the same materials
(Section~\ref{impl:human}); every judgment, human and model alike, is converted into an
eleven-dimensional probability vector (Section~\ref{impl:vectors}); metrics of shift are derived
from those vectors (Section~\ref{impl:metrics}); and a quality-control gate
(Section~\ref{impl:qc}) must be passed before any inferential statistic is calculated
(Section~\ref{impl:stats}).

\section{Prompt Construction and Model Evaluation}
\label{impl:prompts}

\subsection{Stimulus Construction}
\label{impl:stimulus}

\noindent
Each stimulus consists of the unmodified text followed by an optional author cue, in the fixed
format \texttt{"Text: \{text\}"} for the neutral condition and \texttt{"Text: \{text\}\textbackslash
nAuthor: \{cue\}"} for the other three, where \texttt{\{cue\}} is the literal string
\texttt{"a person"} for the generic human condition and the assigned persona description for the
two persona conditions. All four framing conditions keep the text byte-identical; only the cue,
and whether it appears at all, varies.

One property of the source corpus had to be fixed during implementation. crowd-enVent texts were
elicited using the frame \emph{``I felt [emotion] when\slash that\slash if\ldots''}, and in some
texts the author's own emotion word was still present in the stimulus. A text like this gives its
own answer away. In every affected text, the emotion word was redacted and replaced with an
ellipsis, which keeps the sentence structure but removes the answer. Section~\ref{impl:leak}
documents how the issue was detected and corrected after an initial round of model collection.

\subsection{Model Selection}
\label{impl:models}

\noindent
Seven large language models were tested, spanning proprietary, commercial and open-weight
systems: \textbf{OpenAI}, \textbf{Claude}, \textbf{Gemini}, \textbf{Qwen3}, \textbf{Gemma3},
\textbf{Ministral} and \textbf{Llama 3.1}. Proprietary and commercial models were queried through
their respective APIs; open-weight models were run locally through Ollama. The shared collection
script also supported a Moonshot Kimi~K3 endpoint as a fourth proprietary/commercial option, but
no Kimi~K3 output was generated or retained: the corresponding
\texttt{KimiK3\_label}\slash\texttt{KimiK3\_reason} columns are empty across all 7{,}200 rows of the
raw result files, and a dedicated quality-control check confirms Kimi~K3 has no normalised-label
column and is correctly absent from the seven-model set analysed throughout this thesis.

\textbf{Recovered model identifiers.} The defaults set in the collection scripts (not recorded
per-row in the result files themselves, so recovered from the scripts and their run logs rather
than from the data) give: \texttt{gpt-5-nano} (OpenAI), \texttt{claude-haiku-4-5-20251001}
(Claude), \texttt{gemini-3.6-flash} (Gemini), \texttt{qwen3:8b}, \texttt{gemma3:4b},
\texttt{ministral-3:8b} and \texttt{llama3.1:8b} (open-weight, via Ollama). No explicit
temperature, top-$p$, or other decoding parameter was set for any model, so each used its
provider's or Ollama's own default, which is not separately recorded and cannot be recovered.
Collection dates can only be reconstructed partially: run logs of the open-weight models date
their collection to early September 2026 (specifically 2026-09-03 to 2026-09-04); no equivalent
timestamped log was found for the proprietary and commercial models' collection runs. These
identifiers are the defaults actually present in the scripts that produced this data; they are
not independently confirmed against a per-response API version field, since none was retained.

The proposal required three systems; seven were used because of RQ2. What one model shows as a
framing effect could be an artefact of a single training pipeline, but an effect that replicates
across seven independently developed systems is far less plausibly attributed to any single one
of them.

\subsection{Prompt Configurations}
\label{impl:configs}

\noindent
Six prompt configurations were used, formed by crossing three prompt variants with two label
formats:

\begin{itemize}
    \item \textbf{Variants:} zero-shot, one-shot, few-shot.
    \item \textbf{Label formats:} the eleven emotions presented as a numbered list, or embedded
    inline in the question.
\end{itemize}

\noindent
Each model answered every (text, framing) pair once under each configuration, giving
$300 \times 4 \times 3 \times 2 = 7{,}200$ cells per model-ensemble position. Configurations are
\textbf{never pooled at the level of raw responses}. Where an aggregate across configurations is
given, each configuration's results are calculated independently first and then averaged, so that
no analysis treats six correlated observations of one text as six independent observations.

\subsection{Response Normalisation and Off-List Answers}
\label{impl:normalisation}

\noindent
Models were instructed to return exactly one emotion from the fixed set of eleven, and not every
response did. Two kinds of non-compliance were handled explicitly.

\textbf{Formatting variation}, answers that differ only in case, whitespace, or surrounding
punctuation, were normalised to the canonical label.

\textbf{Off-list answers}, where the respondent named an emotion outside the eleven, such as
\emph{anxiety}, \emph{disappointment}, or \emph{frustration}, were \textbf{excluded} rather than
mapped onto the nearest in-set label. Mapping would mean deciding which of the eleven a given
synonym corresponds to, and that decision would itself be an interpretive judgment inserted into
the data by the analyst. Exclusion is the conservative choice: it shrinks the effective ensemble
for the affected cell instead of inventing a vote.

Off-list rates differ substantially by model and are reported in Section~\ref{res:offlist}. The
present analysis does not record any synonym-mapping sensitivity result; whether the exclude-only
policy affects a specific conclusion remains untested here.

\subsection{A Data-Quality Correction Made During Analysis}
\label{impl:leak}

\noindent
Corrections made partway through the study are documented, rather than presenting the corrected
data as though it had always existed. The emotion-word leak described in
Section~\ref{impl:stimulus} was caught and fixed on two separate occasions with different scopes.
Both are described below, not only the more visible second one.

\textbf{First round.} An early audit of the stimulus texts identified 401 candidate texts still
carrying the author's emotion word, of which 102 fell inside the 300-text sample (78 corrected in
place, 24 replaced outright because a clean substitute was not otherwise available). The model
side was fully re-run against the corrected corpus. Since the human side had already been
annotated for 78 of those 102 texts before the fix, those human ratings answered a stimulus that
no longer existed; they were discarded and re-collected, together with the human annotations for
the 24 replacement texts, in a dedicated collection round. The same round also collected the
study's first human ratings for the neutral condition, which up to that point had been supplied
by reusing the original corpus's five-reader labels, a different, disjoint rater pool from the
persona conditions. All of this, 900 human vectors across three framings, each now from data
collected specifically for the post-correction 300-text sample, was in place before the second
round began.

\textbf{Second round.} After an initial round of \emph{model} collection against that corrected
corpus, a further 49 of the 300 texts (16.3\%) were found to still retain the author's emotion word
in the model stimulus, and the same texts additionally carried an outdated persona pair from an
earlier design revision. The corrected wording existed in the master sampling file but had never
been propagated into the model prompt files for that batch. This is a distinct error from the
first round rather than a repetition of it: it affected the \emph{model} prompt files
specifically, was not identified until the human side described above was already complete, and
does not claim to share a common cause with the first round beyond both having occurred within
the same elicitation frame.

The error was corrected through \textbf{full re-collection}, not through exclusion. All 1{,}176
affected rows ($49 \text{ texts} \times 24 \text{ conditions}$) of the stimulus and both personas
were reconstructed from the corrected master file, previous answers for those rows were discarded
since they responded to a prompt that no longer existed, and all seven models were re-queried.

The change in results was material, and that is the reason for reporting it. Model persona
sensitivity rose from 8.7\% to 11.3\% on unambiguous texts and from 12.2\% to 14.8\% on
author-independent texts, while author-relevant texts, of which only nine were affected, shifted
only marginally. The size of this change follows the number of affected texts per category,
consistent with the correction having removed a suppressing artefact. This round,
however, altered two things at once for the affected 49 texts, both the leaked emotion word and
the persona pair itself, so the rise cannot be credited to removing the leak alone. Some or all of
it may instead reflect the updated persona pairs being more effective manipulations than the
outdated ones they replaced. No hypothesis verdict changed qualitatively.

\section{Human Annotation Study}
\label{impl:human}

\noindent
Human annotators saw the same texts, the same eleven-label set and the same author cues as the
models, so both sides were responding to materially identical stimuli.

Each (text, framing) pair was annotated several times by independent annotators, with the aim of
three annotators per cell; some cells had more. As the approved proposal specified,
\textbf{no annotator read the same text more than once}. A participant who rated a text under
one framing never saw that text under a different framing.

The statistical treatment in Section~\ref{impl:stats} rests on this property. Since the
conditions of a given text were rated by disjoint sets of individuals, there is no
individual-level pairing across conditions to preserve. The comparison is therefore
\textbf{between-subjects at the level of framing condition}. Disjoint raters are a documented
property of the collection procedure; exchangeability under a null of no framing effect
additionally requires that the groups assigned to each condition would, absent a framing effect,
have produced the same distribution of judgments, which disjoint recruitment alone does not
guarantee if allocation or the participant pool differed systematically across rounds
(Section~\ref{impl:permutation}).

The neutral condition was collected in a separate, later round from the persona conditions. Only
vote values were stored in the analysis pipeline; respondent identifiers were not retained. The
consequences of this for inference are stated in Section~\ref{impl:stats}.

\textbf{A traceability limit in the neutral collection.} The neutral condition's forms were
initially designed and distributed for one rater per text, and the collection records available
alongside this thesis's data reflect that design. The dataset used throughout this thesis instead
carries three raters per neutral cell, matching the persona conditions. The source of the
additional two raters per text for the neutral condition could not be reconstructed from the
retained collection records at the time of writing. It is disclosed as a limit on this thesis's
own reproducibility, not as an allegation that the three-rater neutral vote counts are invalid:
these are exactly the numbers every downstream calculation in this thesis, including all six H1
tests, is computed from, and nothing in the analysis pipeline is inconsistent with them. But a
reader who wants to trace every human vote back to a raw, timestamped response should know that
this particular step in the collection history is not fully documented in the materials retained
for this thesis.

\section{Emotion Probability Vector Construction}
\label{impl:vectors}

\noindent
The representation described here is the core of the thesis.

\subsection{Definition}
\label{impl:vecdef}

\noindent
Let $E = \langle e_1, \ldots, e_{11}\rangle$ be the fixed, alphabetically ordered emotion set.
For any group of judgments about a single (text, condition), whether those judgments come from
human raters or from models, let $c_i$ be the number of judgments naming emotion $e_i$, and let
$n = \sum_i c_i$ be the number of usable judgments. The emotion probability vector is

\begin{equation}
\mathbf{v} = \left\langle \frac{c_1}{n}, \frac{c_2}{n}, \ldots, \frac{c_{11}}{n} \right\rangle,
\qquad \sum_{i=1}^{11} v_i = 1 .
\end{equation}

\noindent
Every vector in this thesis is constructed this way. A human vector for a (text, framing) pair
aggregates that cell's raters; a model ensemble vector for a (text, framing, variant, format) cell aggregates the seven models' single answers. Off-list and
missing answers are excluded from $c_i$ and reduce $n$ accordingly; they are never silently
counted as zeros, which would distort the denominator.

\subsection{Why Distributions Rather Than Labels}
\label{impl:whyvectors}

\noindent
The reason for this representation is empirical. In the original five-reader corpus, only 80.3\%
of texts have a unique emotion commanding a strict majority. Of the remainder, 6.3\% have a
unique top choice that never reaches 50\% (a plurality, not a majority) and 13.3\% have
\textbf{no unique top choice at all}. For roughly one text in five, a pipeline that reports a
single label must either retain the most frequent label while concealing how divided readers
actually were, or break a tie arbitrarily.

A vector sidesteps the choice. It can also express something a label cannot: the difference
between a judgment that barely moved and one that was substantially reinterpreted but happened to
keep the same most-popular answer. Section~\ref{res:vectorization} measures that difference.

\section{Metrics}
\label{impl:metrics}

\noindent
The vectors give rise to two families of metric: one for the \emph{magnitude} of a shift and one
for its \emph{direction}. They answer different questions and are never conflated.

\subsection{Magnitude: Total Variation Distance}
\label{impl:tv}

\noindent
For two vectors $\mathbf{p}$ and $\mathbf{q}$ over the same emotion set, total variation distance
is

\begin{equation}
\mathrm{TV}(\mathbf{p}, \mathbf{q}) = \frac{1}{2}\sum_{i=1}^{11} |p_i - q_i| \; \in [0, 1].
\end{equation}

\noindent
TV is the share of probability mass that must be moved to turn one distribution into the other.
It is zero when the two are identical and one when their supports are disjoint. Shift magnitudes
are reported as $\mathrm{TV} \times 100$ and referred to as shift percentages.

TV measures magnitude only. It expresses how much moved, but not which emotions moved or in which
direction; the directional metrics below are therefore reported alongside it as a matter of
course, not as an optional extra.

\subsection{Direction: Delta Vectors and Cosine Similarity}
\label{impl:delta}

\noindent
The signed change induced by framing is the delta vector

\begin{equation}
\boldsymbol{\delta} = \mathbf{v}_{\text{persona}} - \mathbf{v}_{\text{neutral}},
\end{equation}

\noindent
whose components sum to zero and record which emotions gained mass and which lost it. To ask
whether a model's shift resembles the human shift \emph{for the same text}, the two delta vectors
are compared by cosine similarity:

\begin{equation}
\cos(\boldsymbol{\delta}_{H}, \boldsymbol{\delta}_{L}) =
\frac{\boldsymbol{\delta}_{H} \cdot \boldsymbol{\delta}_{L}}
     {\lVert \boldsymbol{\delta}_{H}\rVert \, \lVert \boldsymbol{\delta}_{L}\rVert}.
\end{equation}

\noindent
Cosine is undefined when either delta is the zero vector, which occurs whenever one side did not
move at all. These cases are \textbf{counted and reported} rather than silently dropped, because
their frequency is itself a finding (Section~\ref{res:zeroshift}).


\subsection{Dominant Sets and Their Five Outcomes}
\label{impl:dominant}

\noindent
The dominant set of a vector is the set of emotions attaining its maximum probability. It is a
\emph{set} rather than a single label because ties are common at small sample sizes and resolving
them arbitrarily is exactly the practice this thesis is examining.

Comparing dominant sets before and after framing yields five descriptive outcomes:

\begin{itemize}
    \item the sets are \textbf{identical};
    \item a single emotion \textbf{splits into a tie};
    \item a tie \textbf{resolves} to a single emotion;
    \item two ties \textbf{partially overlap};
    \item the sets are \textbf{disjoint}.
\end{itemize}

\noindent
These describe relationships between complete top sets. A selected-label analysis instead
chooses one winner from each set, so its changed/unchanged verdict does not coincide with any one
of these five outcomes; a tie can change the chosen winner while the sets still overlap. The set
outcomes are given in Section~\ref{res:dominant} and the selected-label comparison in
Section~\ref{res:vectorization}.

\subsection{Dispersion Measures}
\label{impl:dispersion}

\noindent
Three further quantities summarise the shape of a vector: Shannon entropy (and its normalised
form), maximum probability as a measure of consensus strength, and support size, the count of
emotions receiving any mass. These address a question distinct from magnitude, namely whether
framing \emph{concentrates} or \emph{disperses} interpretation, and are examined in
Section~\ref{res:entropy}.

\section{Quality-Control Gate}
\label{impl:qc}

\noindent
Every inferential statistic in this thesis was computed only after its inputs passed quality
control. Since the notebook rebuild described in Section~\ref{impl:repro}, this rule is enforced
at two levels rather than through a single external report. First, the data-loading routine that
every notebook calls halts immediately, by assertion, if a structural invariant fails: there is no
code path by which an analysis can run against malformed inputs. Second, an independent
verification script re-derives 30 scientific and execution checks from the raw inputs and confirms
that all three notebooks executed top to bottom with zero errors and consecutive execution counts;
all 30 pass.

The checks fall into four groups:

\begin{itemize}
    \item \textbf{Completeness}: every text present in every derived table; expected row counts
    exactly met ($900$ human vectors yielding $600$ comparisons across two persona sides,
    $3{,}600$ model comparisons
    across two persona sides and six prompt configurations, exactly six configuration rows per
    text and persona side).

    \item \textbf{Mathematical validity}: every probability vector sums to one; every delta
    vector sums to zero; total variation equals the sum of positive signed change; the retained
    vote count for every model comparison is equal on both sides of the comparison.

    \item \textbf{Design integrity}: exactly 100 texts per category; at least three model
    identities common to both sides of every model comparison; no duplicate (text, condition) rows;
    every one of the 300 sampled texts reproduces the classification rule of
    Section~\ref{meth:rule} exactly against the 985-text eligible candidate pool.

    \item \textbf{Source consistency}: every category value maps to a known category, and no
    text is classified inconsistently across its own rows.
\end{itemize}

\noindent
The final group was written to catch a particular bug. The human annotation file disagrees with
itself on nine texts: one category is recorded in the neutral rows and another in the persona
rows, a relic of a reclassification between rounds of collection. Because naive de-duplication
does not remove such a conflict but silently resolves it by whichever row happens to come first,
those nine texts had crept into the human-side analyses under the wrong category, producing an
unbalanced 91/100/109 grouping in place of the designed 100/100/100. The fix is structural: instead of reading the per-row category column of the human annotation file,
the master sampling file is read once, at data-loading time, and the category of every text is
looked up from it. The nine texts are checked explicitly to confirm that the defect is real and has
been correctly resolved, not to certify a list that could otherwise remain stale.

\section{Statistical Procedures}
\label{impl:stats}

\noindent
This section sets out the inferential procedures and, just as important, the unit of analysis for
each test.

\subsection{Unit of Analysis}
\label{impl:unit}

\noindent
The risk that keeps recurring in this dataset is pseudoreplication. A single text provides 2
human comparisons and 12 model comparisons (2 persona sides $\times$ 6 configurations). Treating
those as independent observations inflates $n$ by as much as twelvefold, and correspondingly
inflates significance.

Accordingly, unless stated otherwise, \textbf{the unit of analysis is the text}. Where a statistic
is computed on finer-grained rows, the confidence interval is instead obtained by a bootstrap that resamples
texts rather than rows, so that the correlation among a text's own observations is preserved in
the resampling.

\subsection{Permutation Tests}
\label{impl:permutation}

\noindent
Two different permutation nulls are used \cite{Good2005Permutation}, and the difference between
them is one of substance, not merely of technicality.

\textbf{For the human data (H1)}, the votes are pooled within a text across the two conditions
under comparison, reshuffled, and re-split at the observed group sizes. The different conditions
of a text were rated by different people, so there is no within-text participant pairing
requirement. The permutation also presumes that, given no framing effect, the groups assigned to
that text's conditions would have the same distribution of emotion judgments; disjoint
participants alone do not imply exchangeability if recruitment or allocation differed
systematically.

\textbf{For the model data (H2)}, the same pooling would be invalid. The same seven models answer
both conditions, so model identity is tracked on both sides, and the models are not exchangeable
with each other; Section~\ref{res:h2magnitude} measures directly how differently they behave. The
correct null holds model identity fixed and asks whether each model's own two answers could just
as well have been obtained the other way round. The test is therefore a \textbf{paired sign-flip}: for
each (cell, model) pair independently, the two labels are swapped with probability one half.

Both tests use 10{,}000 iterations with a fixed random seed.

\subsection{Effect Sizes and Confidence Intervals}
\label{impl:effects}

\noindent
Effect sizes are reported alongside every $p$-value. Category contrasts (H3) are reported as the
raw difference in mean shift between groups, with a text-level bootstrap confidence interval,
rather than as a standardised ordinal effect size; matched-pairs rank-biserial correlation is used
for paired comparisons (H4 magnitude), and Spearman's $\rho$ for monotonic association (H4
covariation). Confidence intervals are percentile bootstrap intervals with text-level clustering.

Where an association is summarised by a correlation, inference is drawn from the
\textbf{clustered bootstrap itself}, as opposed to an analytic $p$-value computed over all rows. Over correlated rows, an
analytic $p$ is not merely imprecise but misleading. Duplicating every row would leave $\rho$
unchanged and shrink $p$, the very behaviour clustering is meant to prevent.

\subsection{Multiple Comparisons}
\label{impl:multiple}

\noindent
Two correction regimes are used, and which applies is indicated at each test. Confirmatory
hypothesis families are controlled by \textbf{Holm} correction \cite{Holm1979Sequential}, where
the family-wise error rate should be controlled. Exploratory families, such as the 66 per-emotion
tests (Section~\ref{res:emotions}), where the aim is to limit the proportion of false positives
among discoveries rather than exclude them entirely, are controlled by \textbf{Benjamini--Hochberg}
false discovery rate control \cite{Benjamini1995FDR}.

\subsection{Referencing Against a Sampling Null}
\label{impl:nullref}

\noindent
One last procedure is needed because of an easily overlooked property of small-sample TV.

\textbf{TV between two small vote samples is not zero when there is no effect.} Two independent
three-vote draws from the same underlying distribution will often disagree, and they disagree
more when the underlying distribution is dispersed. Since the ambiguous categories are by
construction the dispersed ones, a raw comparison of TV across categories conflates baseline
dispersion with any framing effect.

To separate the two, each text is given its own reference point. That text's votes are
pooled across the two conditions, reshuffled, and re-split at the observed sizes; the average TV
over many such draws is the TV expected for that specific text if framing did nothing. The
observed value minus this expectation is referred to as the text's \textbf{excess TV}.

This is applied in Section~\ref{res:h3null} as a sensitivity analysis. It is reported next to the
raw contrasts and never replaces them: it was specified after the raw result was known, and
making it the primary analysis would amount to selecting an analysis by its outcome.

\subsection{Reproducibility}
\label{impl:repro}

\noindent
All permutation tests and bootstrap operations are deterministic and were confirmed to reproduce
exactly across separate process invocations with a fixed seed. A single non-determinism bug was
discovered and fixed during development: an unsorted set intersection produced an iteration order
dependent on per-process hash randomisation, causing seeded results to be irreproducible across
runs. All collections of string keys feeding a seeded loop are now explicitly sorted before use.

The analysis notebooks were later rewritten so that the calculations behind the four main
hypotheses are stored directly within the notebook rather than being imported from a separate
module; the core pipeline can now be audited without referring to other files. All 44 result tables
generated by the rewritten notebooks were compared against the results generated prior to the
rewrite and were binary equivalent; the rewrite reformatted, but did not redefine, the
calculations. A small number of additional results added later (the inter-annotator agreement
figures, the conventional majority-label benchmark, the entropy hypothesis tests of
Section~\ref{res:entropy}, and the transition-matrix figure of Section~\ref{res:dominant}) are not
inlined into the notebooks but are generated by independent scripts alongside them.
